\documentclass[12pt]{article}
\usepackage[utf8]{inputenc}
\usepackage[T1]{fontenc}
\usepackage[margin=1in]{geometry}
\usepackage{setspace}
\usepackage{xcolor}
\usepackage[colorlinks=true, urlcolor=blue!60!black, linkcolor=blue!60!black]{hyperref}
\setlength{\parindent}{0em}
\setlength{\parskip}{0.45em}
\renewcommand{\baselinestretch}{1.0}
\begin{document}
\begin{center}
{\Large \textit{Bolling v. Sharpe}, 347 U.S. 497 (1954)\par}
\vspace{0.3em}
{\normalsize Case Brief --- One-Page Edition\par}
\end{center}

\textbf{Memory Jogger.} The D.C. companion to \textit{Brown}: the Fifth Amendment's Due Process Clause bars federal school segregation, launching reverse incorporation (\href{https://supreme.justia.com/cases/federal/us/347/497/}{347 U.S. 497}).

\textbf{Facts.} In 1950 the principal of D.C.'s new Sousa Junior High turned away eleven Black students, including Spottswood Bolling Jr., solely because of their race, despite empty classrooms. Because D.C. is a federal enclave, the Fourteenth Amendment's Equal Protection Clause did not apply.

\textbf{Procedural History.} The district court dismissed; the Court granted certiorari before judgment and decided the case with \textit{Brown} on May 17, 1954.

\textbf{Votes.} Unanimous, 9--0, Chief Justice Warren writing.

\textbf{Holding.} Racial segregation in D.C. public schools violates Fifth Amendment due process. Due process and equal protection are not mutually exclusive, and segregation, unrelated to any proper governmental objective, arbitrarily deprives Black children of liberty.

\textbf{Reasoning.} Both guarantees stem from ``our American ideal of fairness''; racial classifications are ``constitutionally suspect'' under \textit{Korematsu} and \textit{Hi\-ra\-ba\-ya\-shi}; and, given \textit{Brown}, a lesser federal duty would be unthinkable. \textit{Adarand} (1995) later made federal and state equal protection largely congruent.

\textbf{Examples.} \emph{Same-Side:} federal survivor benefits for widows but not widowers fall. \emph{Opposite-Side:} federal alienage limits on Medicare survive rational-basis review. \emph{Fence-Sitter:} excluding Puerto Rico residents from SSI.

\textbf{Critique.} Originalists find no anti-discrimination rule in the 1791 Due Process Clause and see redundancy with the Equal Protection Clause; progressives call the ``unthinkable'' rationale ipse dixit and fault \textit{Adarand}'s rigid congruence.

\textbf{Key Quotation.} ``The `equal protection of the laws' is a more explicit safeguard of prohibited unfairness than `due process of law,' and therefore we do not imply that the two are always interchangeable phrases. But, as this Court has recognized, discrimination may be so unjustifiable as to be violative of due process.'' --- Chief Justice Warren, Opinion of the Court, 347 U.S. at 499.

\end{document}
